% =====================================================================
% Capítulo 1 — Introducción
% =====================================================================
\chapter{Introducción}
\label{cap:introduccion}

\section{Contexto}
\label{sec:contexto}

La calidad de cualquier sistema de adquisición de una señal trabaja
siempre bajo el mismo compromiso: aumenta con el presupuesto invertido
en la medición. Un fotógrafo que dispara de noche alarga la exposición
o sube la sensibilidad del sensor; en todos los casos la moneda de
cambio es la misma —tiempo, energía, exposición— y la factura se paga
en la calidad del resultado.

La tomografía por emisión de positrones (PET, \textit{Positron
Emission Tomography}) es un caso extremo de este compromiso, porque su
presupuesto de adquisición se paga en una moneda especialmente cara:
radiación inyectada al paciente y tiempo inmovilizado dentro de la
máquina. La imagen ---un mapa tridimensional de la actividad metabólica
del cuerpo--- se reconstruye a partir de un \emph{recuento de eventos}
---desintegraciones radiactivas detectadas---, y acumular suficientes
eventos exige una dosis de radiotrazador considerable y entre quince y
cuarenta y cinco minutos de escaneo. Recortar cualquiera de los dos
significa contar menos eventos y, por tanto, degradar la imagen.

En este trabajo se aborda el escenario de \emph{dosis baja}, donde el
escaneo se adquiere con una vigésima parte (1/20) de la dosis estándar:
el ruido es severo y la imagen pierde buena parte de su valor
diagnóstico. Recuperar la imagen de dosis completa a partir de esa
observación degradada es un \emph{problema inverso} mal planteado ---el
ruido destruye información que ninguna operación determinista puede
restituir---, de modo que toda solución pasa por inyectar un
\emph{prior} sobre qué aspecto tiene una imagen PET plausible, y la
calidad de la reconstrucción depende críticamente de la calidad de ese
prior.

Aquí entran los \textbf{modelos de difusión}, modelos generativos
profundos que se entrenan aprendiendo a quitar ruido: se corrompe una
imagen en pasos sucesivos y una red neuronal aprende a revertir cada
paso, de forma que, encadenándolos desde ruido puro, el modelo
sintetiza muestras nuevas de la distribución con la que fue entrenado.
Un modelo de difusión entrenado sobre PET de dosis completa es,
implícitamente, un prior aprendido sobre el espacio de imágenes
anatómicamente plausibles: justo lo que el problema inverso necesita.

Se estudian los dos paradigmas que la literatura reciente plantea para
este problema ---uno \textbf{supervisado}, que condiciona la red con el
escaneo de dosis baja, y otro \textbf{no condicional con guiado}, que
aprende un prior de dosis completa e inyecta la observación en
inferencia---, detallados junto al resto de objetivos en la
\Cref{sec:objetivos}. El trabajo los implementa de extremo a extremo y
los confronta entre sí y contra dos líneas base discriminativas que
cuantifican qué aporta realmente la difusión.

\section{Motivación e interés del tema}
\label{sec:motivacion}

La motivación de este trabajo se articula en tres planos
complementarios: clínico, metodológico y formativo.

En el plano \textbf{clínico}, cualquier técnica que permita reducir
la dosis de radiotrazador o el tiempo en la máquina sin sacrificar
la calidad diagnóstica tiene un impacto directo sobre el paciente.
Una reducción de dosis disminuye la exposición a radiación
ionizante, especialmente relevante en estudios oncológicos de
seguimiento —donde un mismo paciente puede ser escaneado decenas de
veces a lo largo del tratamiento—, en pediatría —donde la
sensibilidad a los efectos de la radiación es mayor— y
en cribados poblacionales, donde el riesgo agregado se reparte
sobre un gran número de individuos sanos. Acortar el tiempo de
escaneo, por su parte, mejora la experiencia del paciente, reduce
los artefactos por movimiento involuntario y aumenta el número de
escaneos diarios que un equipo puede realizar, agilizando así
las listas de espera.

En el plano \textbf{metodológico}, los modelos de difusión son uno de
los paradigmas generativos más activos del momento, y varios trabajos
recientes en imagen médica los sitúan como estado del arte en
reconstrucción de PET de dosis baja, TAC \textit{sparse-view} y
resonancia magnética ultrarrápida~\cite{bjrai2024}. Estudiar a fondo
sus dos formulaciones ---la supervisada y la de guiado posterior--- y
confrontarlas con líneas base discriminativas bien establecidas es una
apuesta razonable por una línea con recorrido.

En el plano \textbf{formativo}, el trabajo encaja directamente con los
contenidos del Máster en Tratamiento Estadístico-Computacional de la
Información: reúne los fundamentos teóricos de los modelos generativos
profundos, el \textit{deep learning} aplicado a visión, la matemática
de los problemas inversos y una capa de ingeniería de sistemas en la
nube. Trabajar con datos sanitarios reales introduce además
restricciones operativas —fugas de información entre particiones,
normalizaciones derivadas únicamente de la observación— ausentes en los
\textit{benchmarks} sintéticos habituales.

\section{Objetivos}
\label{sec:objetivos}

El objetivo general de este Trabajo Fin de Máster es \textbf{diseñar,
implementar y evaluar un sistema de reconstrucción de PET de dosis
baja basado en modelos de difusión, comparando empíricamente dos
paradigmas representativos de la literatura reciente frente a dos
líneas base discriminativas clásicas}. Se concreta en seis objetivos
específicos:

\begin{enumerate}[leftmargin=*]
    \item \textbf{Revisión del estado del arte.} Realizar una
    revisión crítica de la literatura reciente en reconstrucción
    de imagen médica con modelos de difusión, con especial atención
    a las propuestas específicas para PET de dosis baja
    (\Cref{sec:estado-del-arte}). El propósito no es enciclopédico
    sino instrumental: identificar qué paradigmas resultan
    defendibles para el problema concreto y justificar a partir de
    ellos las decisiones de diseño.

    \item \textbf{Apropiación del paradigma de difusión mediante un
    experimento controlado.} Reimplementar de extremo a extremo una
    variante del tutorial canónico de \texttt{diffusers} —un DDPM no
    condicional sobre el conjunto de mariposas del Instituto
    Smithsoniano— ejecutándolo en local. Este experimento preliminar
    valida el flujo completo de entrenamiento e inferencia
    (preprocesamiento, U-Net, \textit{scheduler}, EMA, generación)
    sobre un dominio sencillo antes de abordar el problema principal.

    \item \textbf{Diseño e implementación de dos \textit{pipelines}
    de reconstrucción PET.} Construir, sobre un conjunto de 371
    pares de volúmenes NIfTI (dosis 1/20 / dosis completa), dos
    sistemas de reconstrucción 2D por \textit{slices} axiales: un
    paradigma \textbf{supervisado condicional}, que inyecta el
    escaneo de dosis baja como canal de condicionamiento del U-Net,
    y un paradigma de \textbf{prior no condicional con guiado de
    medición} (\textit{Diffusion Posterior Sampling}), que entrena
    exclusivamente sobre dosis completa e introduce la observación
    en inferencia.

    \item \textbf{Líneas base discriminativas de referencia.}
    Entrenar, con idénticos datos, normalización y presupuesto, dos
    modelos de regresión directa que delimiten la aportación del
    paradigma de difusión: la \emph{misma} U-Net del
    \textit{pipeline} supervisado pero sin proceso de difusión
    (ablación del paradigma generativo) y una RED-CNN, la CNN de
    referencia para \textit{denoising} de imagen médica de dosis
    baja (ablación de la familia arquitectónica). A ellas se suma el
    propio escaneo de dosis baja sin procesar como suelo de
    referencia.

    \item \textbf{Evaluación cuantitativa y cualitativa
    comparativa.} Evaluar los cuatro modelos sobre el mismo conjunto
    de test y los mismos \textit{slices}, con métricas estándar de
    calidad de imagen (PSNR, SSIM, NRMSE) en espacio normalizado,
    métricas de preservación de intensidad en espacio de cuentas, y
    análisis cualitativo sobre volúmenes representativos. La
    comparación debe ser \emph{justa} —mismas particiones, misma
    normalización per-volumen derivada de la dosis baja— para que
    las diferencias sean atribuibles al método y no a artefactos de
    la evaluación.

    \item \textbf{Infraestructura reproducible y análisis de
    costes.} Construir un entorno de ejecución profesional en la
    nube —imagen Docker, dataset preprocesado publicado en Hugging
    Face, GPU bajo demanda en RunPod— que haga cada experimento
    reproducible de extremo a extremo, y documentar explícitamente
    el coste de cada uno de los cuatro entrenamientos, junto con las
    decisiones de ingeniería que condicionan el resultado
    (resolución de trabajo, normalización sin fuga de información,
    elección de \textit{schedule} y de máquina).
\end{enumerate}

En conjunto, estos seis objetivos llevan el trabajo desde la revisión
de la literatura hasta una comparación empírica reproducible y honesta
sobre los recursos que consume.
